\documentclass[11pt,a4paper]{report}

\usepackage[utf8]{inputenc}
\usepackage[T1]{fontenc}
\usepackage[french]{babel}
\usepackage{lmodern}
\usepackage{geometry}
\usepackage{booktabs}
\usepackage{longtable}
\usepackage{array}
\usepackage{tabularx}
\usepackage{graphicx}
\usepackage{xcolor}
\usepackage{hyperref}
\usepackage{listings}
\usepackage{amsmath}
\usepackage{enumitem}
\usepackage{microtype}

\geometry{margin=2.3cm}
\hypersetup{
  colorlinks=true,
  linkcolor=blue!55!black,
  urlcolor=blue!55!black,
  pdftitle={Rapport detaille du projet de classification des adventices},
  pdfauthor={Analyse du projet}
}
\setlength{\parindent}{0pt}
\setlength{\parskip}{0.55em}
\renewcommand{\arraystretch}{1.18}
\setlist[itemize]{leftmargin=1.5em}
\lstset{
  basicstyle=\ttfamily\small,
  breaklines=true,
  frame=single,
  columns=fullflexible
}

\newcommand{\code}[1]{\texttt{\detokenize{#1}}}
\newcommand{\pct}[1]{#1\,\%}

\title{\textbf{Rapport détaillé du projet de classification automatique des adventices}\\
\large Analyse étape par étape des données, méthodes, expériences et résultats}
\author{Projet basé sur la base photographique du Dr Taleb}
\date{Analyse du dépôt au 4 juin 2026}

\begin{document}

\maketitle
\tableofcontents
\listoftables

\chapter*{Résumé exécutif}
\addcontentsline{toc}{chapter}{Résumé exécutif}

Ce projet étudie la reconnaissance automatique de 138 espèces d'adventices à partir
d'une petite base photographique organisée par espèce. Le travail a progressé depuis
l'audit de la base jusqu'à la comparaison de plusieurs familles de méthodes:
Vision Transformers pré-entraînés, transfert par sonde linéaire, évaluation directe
iNaturalist, fine-tuning supervisé, puis BioCLIP et BioCLIP2 en zero-shot et en
apprentissage léger.

L'audit retrouve \textbf{642 images valides} réparties dans \textbf{138 classes},
soit seulement \textbf{4,65 images par classe en moyenne}. Les classes contiennent
entre 2 et 8 images. Le fichier de comptage initial annonce 752 photos; l'écart de
110 éléments correspond exactement à 110 fichiers \code{desktop.ini}, qui ne sont
pas des images. Cette faible quantité de données explique l'intérêt des modèles
pré-entraînés et la difficulté du fine-tuning profond.

Le meilleur résultat global sur les 642 images est obtenu par \textbf{BioCLIP2
zero-shot avec prompts binomiaux}: \textbf{72,59\% de top-1}, \textbf{92,68\% de
top-5} et \textbf{68,90\% de F1 macro}. Dans le protocole plus strict où une image
par classe est réservée au test et où l'expérience est répétée cinq fois, la
meilleure moyenne est de \textbf{71,88\% d'accuracy}, obtenue à la fois par la
sonde logistique BioCLIP2 avec $C=1$ et par le baseline zero-shot à prompts
binomiaux. L'évaluation directe iNaturalist atteint \textbf{89,83\%}, mais uniquement
sur les 59 classes qui peuvent être associées aux étiquettes iNaturalist et lorsque
la décision est restreinte à ces classes; ce résultat n'est donc pas directement
comparable aux expériences sur les 138 classes.

\chapter{Objectif et périmètre du projet}

\section{Problème traité}

Le projet vise à identifier automatiquement l'espèce d'une adventice à partir
d'une photographie. Chaque sous-dossier de la base représente une classe
botanique, dont le nom est généralement un nom scientifique complet, par exemple
\emph{Ammi majus L} ou \emph{Cynodon dactylon (L.) Pers}.

La tâche est difficile pour quatre raisons principales:
\begin{itemize}
  \item le nombre de classes est élevé: 138 espèces;
  \item chaque classe possède très peu d'exemples: 2 à 8 images;
  \item plusieurs espèces sont visuellement proches;
  \item les noms complets comprennent auteurs, parenthèses et variantes
  taxonomiques, alors que les modèles pré-entraînés connaissent souvent seulement
  le binôme genre--espèce.
\end{itemize}

\section{Organisation générale du dépôt}

Le dépôt contient:
\begin{itemize}
  \item la base extraite \code{Bases de données photos Dr Taleb};
  \item une archive source d'environ 1,26 Go;
  \item deux tableaux d'audit: \code{analyse_classes.csv} et
  \code{dataset_audit_counts.csv};
  \item sept scripts Python couvrant les différentes méthodes;
  \item quinze dossiers de résultats expérimentaux;
  \item deux articles PDF conservés comme ressources bibliographiques.
\end{itemize}

\begin{table}[ht]
\centering
\caption{Scripts principaux et rôle}
\begin{tabularx}{\textwidth}{>{\ttfamily\small}p{5.2cm}X}
\toprule
\normalfont Script & Rôle \\
\midrule
evaluate_vit.py & Extraction de caractéristiques ViT et classification par régression logistique. \\
direct_eval_timm_inat21.py & Évaluation directe d'un ViT iNaturalist-2021 sur les classes taxonomiquement associées. \\
finetune_vit_one_test_per_class.py & Fine-tuning supervisé d'un ViT avec une image de test par classe. \\
eval_bioclip_zeroshot.py & Classification zero-shot BioCLIP/BioCLIP2 par similarité image--texte. \\
eval_bioclip_prompt_experiments.py & Comparaison systématique de plusieurs ensembles de prompts. \\
eval_bioclip_linear_probe.py & Sonde logistique et classifieur au centroïde sur les embeddings BioCLIP2. \\
eval_bioclip_advanced_cv.py & Validation croisée répétée, prompts adaptatifs par classe et fusion avec centroïdes. \\
\bottomrule
\end{tabularx}
\end{table}

\chapter{Étape 1: constitution et audit de la base}

\section{Structure de la base}

La base suit le format \emph{ImageFolder}: un dossier racine, puis un sous-dossier
par espèce. Les noms de dossiers servent directement d'étiquettes de classe. Les
scripts utilisent une classe \code{SafeImageFolder} ou une variante équivalente
pour:
\begin{itemize}
  \item accepter les extensions \code{jpg}, \code{jpeg}, \code{png},
  \code{bmp}, \code{webp} et \code{jfif};
  \item convertir les images en RGB;
  \item tolérer les images tronquées;
  \item remplacer une image illisible par une image noire plutôt que d'arrêter
  toute l'expérience.
\end{itemize}

\section{Résultat de l'audit}

\begin{table}[ht]
\centering
\caption{Statistiques consolidées de la base}
\begin{tabular}{lr}
\toprule
Indicateur & Valeur \\
\midrule
Nombre de classes & 138 \\
Nombre annoncé dans \code{analyse_classes.csv} & 752 \\
Images réellement valides & 642 \\
Fichiers \code{desktop.ini} ignorés & 110 \\
Autres fichiers parasites comptés par l'audit & 0 \\
Classes sans différence entre annonce et images & 28 \\
Classes avec un \code{desktop.ini} en plus & 110 \\
Minimum d'images par classe & 2 \\
Maximum d'images par classe & 8 \\
Moyenne d'images par classe & 4,65 \\
\bottomrule
\end{tabular}
\end{table}

Les 642 images valides comprennent 623 fichiers JPEG, 17 PNG et 2 JFIF.
Les classes les moins représentées possèdent deux images, par exemple
\emph{Chenopodium album L}, \emph{Sinapis arvensis L},
\emph{Centaurea eriophora L}, \emph{Bromus rubens L} et
\emph{Cardaria draba (Cav) Desv}. Plusieurs classes atteignent huit images,
notamment \emph{Oxalis pes-caprae L}, \emph{Stellaria media (L.)Vill},
\emph{Urtica urens L}, \emph{Reseda alba L} et \emph{Senecio vulgaris L}.

\section{Conclusion de l'étape}

Le résultat essentiel de cette étape est la mise en évidence d'un problème de
\textbf{classification à très faible nombre d'exemples}. Un apprentissage profond
complet à partir de zéro n'est pas réaliste. La suite du projet s'oriente donc
logiquement vers le transfert d'apprentissage et les modèles vision--langage
pré-entraînés.

\chapter{Étape 2: première référence ViT avec sonde linéaire}

\section{Méthode}

Le script \code{evaluate_vit.py} utilise un modèle pré-entraîné comme extracteur
figé de caractéristiques. Deux backends sont prévus: \code{timm} et Hugging Face.
Les images sont transformées avec le prétraitement propre au modèle, puis les
vecteurs extraits sont classifiés par une régression logistique équilibrée.

La séparation utilise \pct{30} des images pour le test, avec la graine 42. Elle
est stratifiée car chaque classe contient au moins deux images. La régression
logistique utilise $C=1$, un maximum de 2000 itérations et
\code{class_weight=balanced}.

\section{Expériences et résultats}

\begin{table}[ht]
\centering
\caption{Résultats des sondes linéaires ViT}
\begin{tabular}{p{5.2cm}rrrr}
\toprule
Backbone & Accuracy & Acc. équilibrée & F1 macro & F1 pondéré \\
\midrule
ViT Base AugReg IN21K puis IN1K & \pct{52,85} & \pct{50,00} & \pct{45,21} & \pct{48,55} \\
ViT Base pré-entraîné iNaturalist & \pct{32,12} & \pct{30,80} & \pct{26,48} & \pct{27,77} \\
\bottomrule
\end{tabular}
\end{table}

\section{Interprétation}

Le ViT généraliste pré-entraîné sur ImageNet obtient ici une meilleure sonde
linéaire que le ViT Hugging Face spécialisé iNaturalist. Cela ne signifie pas
nécessairement qu'ImageNet est plus pertinent: les couches, le prétraitement et
la qualité des embeddings utilisés comme représentation figée peuvent différer.
Le résultat de \pct{52,85} constitue une première référence supervisée, mais reste
nettement inférieur aux meilleurs résultats BioCLIP2 obtenus ensuite.

\chapter{Étape 3: évaluation directe d'un modèle iNaturalist-2021}

\section{Méthode}

Le script \code{direct_eval_timm_inat21.py} charge le modèle
\code{vit_large_patch14_clip_336.laion2b_ft_augreg_inat21}. Ce modèle possède
10\,000 étiquettes iNaturalist-2021. Le script normalise les noms, extrait le
binôme genre--espèce et construit une correspondance entre les noms du dataset
et les noms iNaturalist.

Deux décisions sont calculées:
\begin{itemize}
  \item une prédiction parmi les 10\,000 étiquettes originales;
  \item une prédiction restreinte aux classes iNaturalist effectivement associées
  à la base locale.
\end{itemize}

Une image par classe associée est sélectionnée avec la graine 42.

\section{Résultat}

\begin{table}[ht]
\centering
\caption{Évaluation directe iNaturalist-2021}
\begin{tabular}{lr}
\toprule
Indicateur & Valeur \\
\midrule
Classes locales totales & 138 \\
Classes associées à iNaturalist-2021 & 59 \\
Classes absentes/non associées & 79 \\
Images évaluées & 59 \\
Accuracy parmi les 10\,000 étiquettes & \pct{74,58} \\
Accuracy restreinte aux 59 classes associées & \pct{89,83} \\
\bottomrule
\end{tabular}
\end{table}

\section{Interprétation}

Le résultat démontre la forte valeur d'un modèle déjà entraîné à reconnaître des
taxons. Cependant, il couvre seulement 59 classes sur 138. L'accuracy restreinte
de \pct{89,83} est excellente mais bénéficie d'un espace de décision fortement
réduit; elle ne doit pas être présentée comme une accuracy sur la base complète.
Le fichier \code{missing_classes.txt} documente les 79 classes non couvertes.

\chapter{Étape 4: fine-tuning supervisé du ViT iNaturalist}

\section{Protocole}

Le script \code{finetune_vit_one_test_per_class.py} réserve exactement une image
par classe au test. Le découpage contient donc:
\[
N_{\text{test}}=138,\qquad N_{\text{train}}=642-138=504.
\]

Les augmentations d'entraînement comprennent un recadrage redimensionné aléatoire,
un retournement horizontal, une variation de couleur et éventuellement un
\emph{random erasing}. L'optimiseur est AdamW avec un scheduler cosinus. Le modèle
est évalué à chaque epoch et l'état ayant le meilleur F1 macro est sauvegardé.

\section{Variantes exécutées}

\begin{table}[ht]
\centering
\caption{Configurations de fine-tuning}
\resizebox{\textwidth}{!}{%
\begin{tabular}{lrrrrrrr}
\toprule
Variante & LR & WD & Lissage & Dropout & Attn. drop. & Blocs libres & Effacement \\
\midrule
Initiale & $2\cdot10^{-5}$ & 0,05 & non enregistré & non enregistré & non enregistré & non enregistré & non enregistré \\
Régularisée forte & $10^{-5}$ & 0,10 & 0,10 & 0,15 & 0,10 & 2 & non enregistré \\
Régularisée légère & $2\cdot10^{-5}$ & 0,10 & 0,05 & 0,05 & 0,00 & 12 & 0,10 \\
\bottomrule
\end{tabular}}
\end{table}

Toutes les variantes utilisent 12 epochs, des images $224\times224$, un batch de
8 et la graine 42.

\section{Résultats}

\begin{table}[ht]
\centering
\caption{Résultats du fine-tuning ViT}
\begin{tabular}{lrrr}
\toprule
Variante & Meilleur epoch & Accuracy & F1 macro \\
\midrule
Initiale & 10 & \pct{20,29} & \pct{14,65} \\
Régularisée forte, 2 derniers blocs & 7 & \pct{2,17} & \pct{1,04} \\
Régularisée légère, 12 blocs & 10 & \pct{20,29} & \pct{14,45} \\
\bottomrule
\end{tabular}
\end{table}

\section{Interprétation}

Le fine-tuning est beaucoup moins performant que l'utilisation de représentations
figées. Avec seulement 504 images d'entraînement pour 138 classes, le modèle ne
dispose en moyenne que de 3,65 images d'entraînement par classe. La variante
fortement régularisée et presque entièrement gelée sous-apprend clairement.
Libérer davantage de blocs remonte l'accuracy à \pct{20,29}, mais le résultat reste
faible.

Il existe aussi une limite méthodologique: le jeu de test est évalué à chaque
epoch et sert à choisir le meilleur état. Il joue donc partiellement le rôle d'un
jeu de validation. Une évaluation finale rigoureuse nécessiterait un troisième
jeu indépendant ou une validation croisée imbriquée.

\chapter{Étape 5: passage à BioCLIP puis BioCLIP2 en zero-shot}

\section{Principe}

BioCLIP associe images et textes dans un même espace vectoriel. Pour une image
$x$ et une classe décrite par un texte $t_c$, le score est calculé par similarité:
\[
s_c = 100\,\frac{f_I(x)}{\lVert f_I(x)\rVert}
\cdot
\frac{f_T(t_c)}{\lVert f_T(t_c)\rVert}.
\]
La classe prédite est celle dont le score est maximal. Aucun classifieur n'est
entraîné sur les images locales dans le mode zero-shot.

\section{Première expérience BioCLIP}

Avec BioCLIP première version, cinq prompts fondés sur le nom complet sont moyennés.
Une image par classe est évaluée.

\begin{table}[ht]
\centering
\caption{Comparaison BioCLIP et BioCLIP2 sur une image par classe}
\begin{tabular}{lrrrr}
\toprule
Modèle & Top-1 & Top-5 & Accuracy équilibrée & F1 macro \\
\midrule
BioCLIP & \pct{36,96} & \pct{64,49} & \pct{36,96} & \pct{28,56} \\
BioCLIP2 & \pct{65,94} & \pct{89,86} & \pct{65,94} & \pct{59,22} \\
\bottomrule
\end{tabular}
\end{table}

Le passage à BioCLIP2 produit un gain majeur de près de 29 points de top-1 sur le
même type de protocole.

\section{Évaluation BioCLIP2 sur toute la base}

Avec les prompts fondés sur le nom complet et les 642 images, BioCLIP2 atteint:
\begin{itemize}
  \item top-1: \pct{68,38};
  \item top-5: \pct{91,59};
  \item accuracy équilibrée: \pct{67,03};
  \item F1 macro: \pct{63,90}.
\end{itemize}

\chapter{Étape 6: étude systématique des prompts BioCLIP2}

\section{Motivation}

Les noms de dossiers contiennent souvent les auteurs botaniques, alors que
l'information la plus discriminante et la plus standard est généralement le
binôme genre--espèce. Le script \code{eval_bioclip_prompt_experiments.py} compare
donc cinq stratégies sur les 642 images.

\begin{table}[ht]
\centering
\caption{Résultats des stratégies de prompts BioCLIP2}
\resizebox{\textwidth}{!}{%
\begin{tabular}{lrrrr}
\toprule
Jeu de prompts & Top-1 & Top-5 & Accuracy équilibrée & F1 macro \\
\midrule
Nom complet, prompt simple & \pct{68,38} & \pct{90,65} & \pct{66,92} & \pct{62,62} \\
Nom complet, cinq prompts plante & \pct{68,38} & \pct{91,59} & \pct{67,03} & \pct{63,90} \\
Binôme, prompt simple & \pct{71,18} & \pct{92,37} & \pct{70,27} & \pct{66,66} \\
Binôme, cinq prompts plante & \textbf{\pct{72,59}} & \textbf{\pct{92,68}} & \textbf{\pct{71,52}} & \textbf{\pct{68,90}} \\
Ensemble nom complet + binôme & \pct{72,12} & \pct{92,06} & \pct{71,02} & \pct{67,96} \\
\bottomrule
\end{tabular}}
\end{table}

\section{Résultat de l'étape}

Le meilleur jeu est \code{binomial_plant_templates}. Il retire les auteurs
botaniques et utilise cinq formulations:
\begin{itemize}
  \item photo du binôme;
  \item photo de la plante;
  \item photo de terrain de l'espèce adventice;
  \item photo botanique;
  \item gros plan des feuilles ou fleurs.
\end{itemize}

Cette modification fait progresser le top-1 de \pct{68,38} à \pct{72,59}.
Le résultat confirme que la normalisation taxonomique des étiquettes est une
partie déterminante du système.

\section{Analyse par classe}

Avec les prompts binomiaux, plusieurs classes obtiennent un F1 de 1, notamment
\emph{Borago officinalis L}, \emph{Heliotropium europaeum L},
\emph{Lamium amplexicaule L}, \emph{Fallopia convolvulus (L.) A. Löve} et
\emph{Senecio vulgaris L}. À l'inverse, certaines classes ont un F1 nul, par
exemple \emph{Chamaemelum mixtum (L.) All}, \emph{Lavatera trimestris L},
\emph{Lathyrus cicera L}, \emph{Geropogon hybridus (L.) Sch.Bip} et
\emph{Aster squamatus (Spreng.) Hieron}. Ces classes constituent des priorités
pour l'enrichissement de la base et l'analyse des confusions.

\chapter{Étape 7: sonde linéaire et centroïdes sur BioCLIP2}

\section{Méthode}

Le script \code{eval_bioclip_linear_probe.py} extrait une seule fois les embeddings
BioCLIP2 des 642 images. Pour chacune des graines 42 à 46, une image par classe est
réservée au test. Le protocole contient donc 504 images d'entraînement et 138
images de test à chaque répétition.

Deux classifieurs sont comparés:
\begin{itemize}
  \item un pipeline \code{StandardScaler} + régression logistique équilibrée,
  avec $C\in\{0{,}01;0{,}1;1;10\}$;
  \item un classifieur au centroïde le plus proche, utilisant la similarité cosinus.
\end{itemize}

\section{Résultats répétés}

\begin{table}[ht]
\centering
\caption{Sondes BioCLIP2, moyenne et écart-type sur cinq découpages}
\begin{tabular}{lrrr}
\toprule
Méthode & Accuracy moyenne & Écart-type & F1 macro moyen \\
\midrule
Logistique, $C=0{,}01$ & \pct{70,58} & \pct{2,77} & \pct{65,08} \\
Logistique, $C=0{,}1$ & \pct{71,01} & \pct{2,47} & \pct{65,17} \\
Logistique, $C=1$ & \textbf{\pct{71,88}} & \pct{2,69} & \textbf{\pct{66,17}} \\
Logistique, $C=10$ & \pct{70,43} & \pct{2,53} & \pct{64,64} \\
Centroïde le plus proche & \pct{63,33} & \pct{2,36} & \pct{57,71} \\
\bottomrule
\end{tabular}
\end{table}

Pour $C=1$, les cinq accuracies individuelles sont \pct{75,36}, \pct{68,12},
\pct{69,57}, \pct{73,91} et \pct{72,46}. La répétition montre que la performance
dépend sensiblement du choix de l'unique image de test de chaque classe.

\section{Conclusion de l'étape}

La sonde logistique BioCLIP2 atteint la même accuracy moyenne que le meilleur
baseline zero-shot répété, avec un F1 macro légèrement supérieur. Elle est simple,
rapide et mieux adaptée à la petite base qu'un fine-tuning profond.

\chapter{Étape 8: validation croisée avancée et prompts adaptatifs}

\section{Méthodes testées}

Le script \code{eval_bioclip_advanced_cv.py} reprend les cinq découpages précédents
et compare:
\begin{enumerate}
  \item le baseline de cinq prompts binomiaux moyennés;
  \item la sélection, pour chaque classe, du meilleur prompt parmi 13 candidats
  à partir des images d'entraînement;
  \item une fusion entre le prompt adaptatif et le centroïde visuel de la classe.
\end{enumerate}

Pour la fusion, un coefficient $\alpha$ est choisi parmi 21 valeurs de 0 à 1:
\[
p_{\text{fusion}} =
\alpha\,p_{\text{texte}}+(1-\alpha)\,p_{\text{centroïde}}.
\]
Les scores de centroïdes utilisés pour choisir $\alpha$ sur l'entraînement sont
calculés en leave-one-out afin de réduire l'optimisme.

\section{Résultats}

\begin{table}[ht]
\centering
\caption{Validation croisée avancée BioCLIP2}
\begin{tabular}{lrrr}
\toprule
Méthode & Accuracy moyenne & Écart-type & F1 macro moyen \\
\midrule
Baseline prompts binomiaux & \textbf{\pct{71,88}} & \pct{3,59} & \textbf{\pct{65,89}} \\
Prompts adaptatifs par classe & \pct{68,12} & \pct{3,34} & \pct{61,01} \\
Prompts adaptatifs + centroïdes & \pct{70,14} & \pct{2,65} & \pct{63,48} \\
\bottomrule
\end{tabular}
\end{table}

Les coefficients texte choisis pour la fusion sont 0,55; 0,85; 0,55; 0,65 et
0,55 selon les cinq découpages.

\section{Interprétation}

La sélection adaptative n'améliore pas le baseline. Avec seulement quelques
images d'entraînement par classe, choisir un prompt distinct pour chaque espèce
sur-apprend facilement les exemples disponibles. La fusion avec les centroïdes
réduit l'écart-type et récupère une partie de la performance, mais reste en dessous
du baseline binomial simple. Cette expérience négative est utile: elle montre que
la simplicité est ici plus robuste.

\chapter{Synthèse comparative}

\begin{longtable}{p{5.9cm}p{3.2cm}rr}
\caption{Vue d'ensemble des principaux résultats}\label{tab:synthese}\\
\toprule
Méthode & Protocole & Accuracy/top-1 & F1 macro \\
\midrule
\endfirsthead
\toprule
Méthode & Protocole & Accuracy/top-1 & F1 macro \\
\midrule
\endhead
ViT AugReg + sonde logistique & split test \pct{30} & \pct{52,85} & \pct{45,21} \\
ViT iNaturalist + sonde logistique & split test \pct{30} & \pct{32,12} & \pct{26,48} \\
ViT iNat21 direct, 10\,000 sorties & 59 classes associées & \pct{74,58} & -- \\
ViT iNat21 direct, sorties restreintes & 59 classes associées & \pct{89,83} & -- \\
Fine-tuning ViT initial & 1 test/classe & \pct{20,29} & \pct{14,65} \\
BioCLIP zero-shot & 1 test/classe & \pct{36,96} & \pct{28,56} \\
BioCLIP2 zero-shot, noms complets & 1 test/classe & \pct{65,94} & \pct{59,22} \\
BioCLIP2 zero-shot, noms complets & toutes les images & \pct{68,38} & \pct{63,90} \\
BioCLIP2 zero-shot, binômes & toutes les images & \textbf{\pct{72,59}} & \textbf{\pct{68,90}} \\
BioCLIP2 sonde logistique $C=1$ & 5 fois 1 test/classe & \pct{71,88} & \pct{66,17} \\
BioCLIP2 baseline binomial & 5 fois 1 test/classe & \pct{71,88} & \pct{65,89} \\
\bottomrule
\end{longtable}

\section{Réponse principale du projet}

Pour les 138 espèces et avec la quantité actuelle de données, la stratégie la plus
efficace est d'exploiter \textbf{BioCLIP2} et de décrire les classes par leur
\textbf{binôme taxonomique}. Le fine-tuning profond n'est pas adapté au volume
actuel. Une sonde logistique sur embeddings BioCLIP2 est compétitive dans un
protocole supervisé; le zero-shot binomial reste particulièrement intéressant car
il ne nécessite aucun entraînement local.

\chapter{Productions et traçabilité}

\section{Fichiers produits par expérience}

Les scripts sauvegardent systématiquement les paramètres, métriques et détails
nécessaires à l'analyse:
\begin{itemize}
  \item \code{metrics.json}: paramètres et métriques globales;
  \item \code{classification_report.json}: précision, rappel et F1 par classe;
  \item \code{confusion_matrix.csv}: matrice de confusion;
  \item \code{predictions.csv}: prédiction individuelle de chaque image;
  \item \code{classes.txt}: correspondance indice--classe;
  \item \code{split_train.csv}, \code{split_test_one_per_class.csv} ou
  \code{split_eval.csv}: traçabilité des découpages;
  \item fichiers \code{*_runs.csv} et \code{*_summary.csv}: résultats détaillés
  et agrégés des expériences répétées;
  \item dossiers \code{best_model}: poids et configuration des ViT fine-tunés.
\end{itemize}

\section{Chronologie expérimentale reconstruite}

D'après les scripts, les paramètres enregistrés et les dates des dossiers de
résultats, la démarche réalisée est la suivante:
\begin{enumerate}
  \item préparation et audit de la base photographique;
  \item création d'une référence ViT généraliste avec sonde linéaire;
  \item test d'un ViT Hugging Face spécialisé iNaturalist;
  \item évaluation directe d'un grand ViT iNaturalist-2021;
  \item essais successifs de fine-tuning, puis ajout de régularisation;
  \item première évaluation BioCLIP en zero-shot;
  \item passage à BioCLIP2, avec fort gain de performance;
  \item évaluation BioCLIP2 sur toute la base;
  \item comparaison de plusieurs formulations textuelles;
  \item adoption des noms binomiaux comme meilleure stratégie de prompts;
  \item apprentissage de sondes logistiques BioCLIP2 sur cinq découpages;
  \item comparaison à un classifieur par centroïdes;
  \item validation croisée avancée avec prompts adaptatifs et fusion multimodale.
\end{enumerate}

\chapter{Limites et points de vigilance}

\section{Taille et diversité de la base}

La limite dominante reste la quantité de données. Deux à huit images par classe
ne suffisent pas pour mesurer correctement la variabilité réelle d'une espèce:
stade de croissance, saison, organe photographié, éclairage, arrière-plan et
appareil. Les résultats peuvent aussi être optimistes si plusieurs images sont
très proches ou proviennent de la même plante.

\section{Comparabilité des protocoles}

Les résultats ne sont pas tous directement comparables:
\begin{itemize}
  \item le split ViT utilise \pct{30} de test;
  \item plusieurs expériences utilisent exactement une image de test par classe;
  \item les résultats zero-shot sur toute la base évaluent les 642 images;
  \item l'évaluation directe iNaturalist couvre seulement 59 classes;
  \item les expériences répétées donnent une moyenne plus fiable qu'un unique split.
\end{itemize}

\section{Validation et sélection de modèles}

Le choix du meilleur epoch de fine-tuning à partir du jeu de test introduit une
fuite de validation. De même, les valeurs de $C$ de la sonde logistique sont
comparées sur les mêmes cinq jeux de test; une validation imbriquée serait plus
rigoureuse pour annoncer une performance finale sans biais de sélection.

\section{Qualité taxonomique}

Le projet montre que les chaînes de caractères des noms d'espèces influencent
fortement BioCLIP2 et le mapping iNaturalist. Des fautes de frappe, synonymes ou
noms anciens peuvent diminuer la performance. Une validation taxonomique externe
des 138 noms et l'utilisation d'identifiants taxonomiques stables amélioreraient
la robustesse.

\section{Mesures manquantes}

Le dépôt ne contient pas de mesure indépendante sur des photos de terrain
entièrement nouvelles, ni d'analyse explicite des doublons visuels. Il ne contient
pas non plus de fichier d'environnement fixant précisément les versions de Python,
PyTorch, Transformers, timm, OpenCLIP et scikit-learn. Ces éléments sont importants
pour reproduire exactement les résultats.

\chapter{Recommandations pour la suite}

\begin{enumerate}
  \item Enrichir chaque classe jusqu'à plusieurs dizaines d'images réellement
  indépendantes, avec différents stades et conditions de prise de vue.
  \item Vérifier et normaliser les 138 noms taxonomiques, puis conserver pour
  chaque classe le nom original, le binôme accepté et un identifiant taxonomique.
  \item Détecter les doublons et groupes de photos d'une même plante, puis séparer
  les groupes plutôt que les images entre entraînement et test.
  \item Conserver BioCLIP2 à prompts binomiaux comme baseline principale.
  \item Évaluer la sonde logistique BioCLIP2 avec validation croisée imbriquée.
  \item Réserver un jeu de test final jamais utilisé pendant le développement.
  \item Tester un regroupement hiérarchique famille--genre--espèce et utiliser le
  top-5, déjà supérieur à \pct{92}, dans une interface d'aide à l'identification.
  \item Ajouter un fichier \code{requirements.txt} ou un environnement Conda et
  documenter les commandes exactes de reproduction.
\end{enumerate}

\chapter{Conclusion générale}

Le projet a exploré de façon progressive et cohérente plusieurs solutions de
classification botanique dans un contexte de données extrêmement limitées.
Les premières expériences ViT ont établi des références, l'évaluation iNaturalist
a montré l'intérêt de la connaissance taxonomique pré-apprise, et les essais de
fine-tuning ont démontré les limites de l'apprentissage profond supervisé sur
seulement quelques images par classe.

La contribution expérimentale la plus importante est l'identification de BioCLIP2
et des prompts binomiaux comme meilleure combinaison actuelle. Elle atteint
\pct{72,59} de top-1 et \pct{92,68} de top-5 sur l'ensemble des 642 images.
Les expériences répétées confirment une accuracy moyenne proche de \pct{72}.
Le projet produit également une traçabilité riche grâce aux prédictions,
matrices de confusion, rapports par classe et résumés de validation croisée.

La prochaine amélioration décisive ne viendra probablement pas d'un modèle plus
complexe, mais d'une base plus grande, mieux normalisée taxonomiquement et évaluée
sur un jeu externe indépendant.

\appendix
\chapter{Inventaire détaillé des dossiers de résultats}

\begin{longtable}{p{7.4cm}p{7.1cm}}
\toprule
Dossier & Contenu/interprétation \\
\midrule
\endfirsthead
\toprule
Dossier & Contenu/interprétation \\
\midrule
\endhead
\code{results_vit_eval} & Sonde logistique sur ViT généraliste timm. \\
\code{results_vit_hf_inaturalist_eval} & Sonde logistique sur ViT Hugging Face iNaturalist. \\
\code{results_vit_inat21_direct_one_test_per_matched_class} & Évaluation directe sur les 59 classes associées à iNat21. \\
\code{results_vit_finetuned_one_test_per_class} & Premier fine-tuning ViT. \\
\code{results_vit_finetuned_regularized_one_test_per_class} & Fine-tuning fortement régularisé, deux blocs libres. \\
\code{results_vit_finetuned_mild_regularized_one_test_per_class} & Fine-tuning légèrement régularisé, douze blocs libres. \\
\code{results_bioclip_zeroshot_one_test_per_class} & BioCLIP v1 zero-shot. \\
\code{results_bioclip2_zeroshot_one_test_per_class} & BioCLIP2 zero-shot, une image par classe. \\
\code{results_bioclip2_zeroshot_all_dataset} & BioCLIP2 zero-shot sur les 642 images, noms complets. \\
\code{results_bioclip2_zeroshot_binomial_prompts_all_dataset} & BioCLIP2 zero-shot sur les 642 images, prompts binomiaux. \\
\code{results_bioclip2_prompt_experiments_all_dataset} & Comparaison de cinq jeux de prompts. \\
\code{results_bioclip2_linear_probe_repeated_one_test_per_class} & Sondes logistiques répétées. \\
\code{results_bioclip2_linear_probe_and_centroid_repeated_one_test_per_class} & Sondes logistiques et centroïdes répétés. \\
\code{results_bioclip2_advanced_cv} & Prompts adaptatifs, centroïdes et validation croisée avancée. \\
\bottomrule
\end{longtable}

\chapter{Formules des métriques}

Pour $K$ classes, l'accuracy mesure la proportion totale de prédictions correctes:
\[
\mathrm{Accuracy}=\frac{\text{nombre de prédictions correctes}}{\text{nombre total d'exemples}}.
\]

L'accuracy équilibrée est la moyenne du rappel des classes:
\[
\mathrm{BalancedAccuracy}=\frac{1}{K}\sum_{c=1}^{K}\mathrm{Recall}_c.
\]

Le F1 d'une classe combine précision et rappel:
\[
F1_c=2\frac{\mathrm{Precision}_c\,\mathrm{Recall}_c}
{\mathrm{Precision}_c+\mathrm{Recall}_c}.
\]

Le F1 macro est la moyenne non pondérée des F1 de toutes les classes:
\[
F1_{\mathrm{macro}}=\frac{1}{K}\sum_{c=1}^{K}F1_c.
\]
Il est particulièrement pertinent ici car chaque espèce doit compter de manière
égale malgré le léger déséquilibre du nombre d'images.

\chapter{Ressources présentes dans le dépôt}

Deux articles PDF sont présents à la racine:
\begin{itemize}
  \item \code{paper-178.pdf};
  \item \code{Zhou_BBN_Bilateral-Branch_Network_With_Cumulative_Learning_for_Long-Tailed_Visual_Recognition_CVPR_2020_paper.pdf}.
\end{itemize}

Le second concerne explicitement la reconnaissance visuelle à longue traîne.
Ces documents constituent des ressources bibliographiques pertinentes, mais aucun
appel direct à une implémentation BBN n'apparaît dans les scripts analysés.

\end{document}
